\documentclass{beamer}

\usepackage[francais]{babel}
\usepackage{rtxslides}
\usepackage{tikz}

\title{\rtx\ -- Technical Assessment 2}
\author{Louis Opter \\ \texttt{<louis@lse.epitech.eu>}}

\begin{document}

\begin{frame}
\titlepage
\end{frame}

\begin{frame}
\rmfamily{\itshape{«~\rtx\ est un \textcolor{rathaxesred}{langage dédié} qui
permet de \textcolor{rathaxesred}{décrire des pilotes} de périphériques. \rtx\
compile vers des modules noyaux écrits en C pour \textcolor{rathaxesred}{Linux},
\textcolor{rathaxesred}{Windows} et \textcolor{rathaxesred}{OpenBSD}.~»}}
\end{frame}

\begin{frame}{Développement de pilotes}
\begin{itemize}
\item code critique ;
\item double compétence ;
\item code spécifique ;
\item auto formation.
\end{itemize}
\end{frame}

\begin{frame}{Objectifs de Rathaxes}
\begin{itemize}
\item imposer plus de vérifications ;
\item séparer les compétences ;
\item unifier le développement.
\end{itemize}
\end{frame}

\begin{frame}{Cas d'utilisations}

\begin{center}
\only<1>{\LARGE\itshape{{Trois cas d'utilisations :}}}

\only<2>{\begin{tikzpicture}
\draw (-4, 1.5) node[above, text width=2cm, text centered, rotate=45]{\small{Développeur de pilotes}} ;
\draw (-4, -1) node[below]{\includegraphics[width=2cm]{logos/nvidia_white}} ;
\draw (-4, 0) node[below]{\includegraphics[width=2cm]{logos/intel}} ;
\draw (-4, 0.5) node[below]{\includegraphics[width=2cm]{logos/broadcom_white}} ;
\draw (-4, -3) node{\ldots} ;

\draw (0, 1.5) node[above, text width=2cm, text centered, rotate=45]{\textcolor{black}{\small{Mainteneur Rathaxes}}} ;

\draw (4, 1.5) node[above, text width=2cm, text centered, rotate=45]{\textcolor{black}{\small{Développeur systèmes}}} ;
\fill[black] (4, 0) circle(1.13) ; % hidden placeholder to make the right overflow
\end{tikzpicture}}

\only<3>{\begin{tikzpicture}
\draw (-4, 1.5) node[above, text width=2cm, text centered, rotate=45]{\small{Développeur de pilotes}} ;
\draw (-4, -1) node[below]{\includegraphics[width=2cm]{logos/nvidia_white}} ;
\draw (-4, 0) node[below]{\includegraphics[width=2cm]{logos/intel}} ;
\draw (-4, 0.5) node[below]{\includegraphics[width=2cm]{logos/broadcom_white}} ;
\draw (-4, -3) node{\ldots} ;

\draw (0, 1.5) node[above, text width=2cm, text centered, rotate=45]{\textcolor{black}{\small{Mainteneur Rathaxes}}} ;

\draw (4, 1.5) node[above, text width=2cm, text centered, rotate=45]{\small{Développeur systèmes}} ;
\draw (4, 0.5) node[below]{\includegraphics[width=2cm]{logos/tux}} ;
\draw (4, 0) node[above]{\includegraphics[width=2cm]{logos/windows_7_white}} ;
\draw (4, -3) node{\ldots} ;
\end{tikzpicture}}

\only<4>{\begin{tikzpicture}
\draw (-4, 1.5) node[above, text width=2cm, text centered, rotate=45]{\small{Développeur de pilotes}} ;
\draw (-4, -1) node[below]{\includegraphics[width=2cm]{logos/nvidia_white}} ;
\draw (-4, 0) node[below]{\includegraphics[width=2cm]{logos/intel}} ;
\draw (-4, 0.5) node[below]{\includegraphics[width=2cm]{logos/broadcom_white}} ;
\draw (-4, -3) node{\ldots} ;

\draw (0, 1.5) node[above, text width=2cm, text centered, rotate=45]{\small{Mainteneur Rathaxes}} ;
\draw (-0.5, 0.5) node[below]{\includegraphics[width=2cm]{logo_latex}} ;

\draw (4, 1.5) node[above, text width=2cm, text centered, rotate=45]{\small{Développeur systèmes}} ;
\draw (4, 0.5) node[below]{\includegraphics[width=2cm]{logos/tux}} ;
\draw (4, 0) node[above]{\includegraphics[width=2cm]{logos/windows_7_white}} ;
\draw (4, -3) node{\ldots} ;
\end{tikzpicture}}
\end{center}

% Note: I just discovered that beamer support columns inside frame, I'll give a
% try next time…

\transdissolve<2>
\transdissolve<3>
\transdissolve<4>
\end{frame}

\begin{frame}{Architecture}
\begin{center}
\only<1>{\LARGE{\itshape{Séparée en plusieurs composants :}}}

\only<2>{\LARGE{\itshape{\textasciitilde~Pause Valise~\textasciitilde}}}
\end{center}
\end{frame}

\begin{frame}{Merci}
\begin{center}
\Large{\itshape{Questions ?}}
\end{center}
\end{frame}

\end{document}
